\section{The first correction at the corner}\label{sec:asym}

The reference square determines the leading terms and the constant,
but it does not determine the first correction.  For that we return to
the exact row decomposition \eqref{eq:corner-exact-decomposition}.
There are three contributions to track: the scalar $-1+2/m$, the
harmonic triangle $\mathcal H_m$, and the smooth residual $4V_n/(nm)$.
Each has a first correction.  Their combination is simpler than any
one of them.

Write
\begin{equation}\label{eq:asym-normalization}
 \begin{gathered}
 m=n+1,\qquad N=n+\tfrac12=m-\tfrac12,\\
 a=\frac2\pi,\qquad d=\gamma+\log\frac8\pi,\qquad
 A(z)=a(\log z+d).
 \end{gathered}
\end{equation}
The constants used in the expanded form of the answer are
\begin{equation}\label{eq:asym-coefficients}
 c_1=ad,\qquad b=2ac_1=\frac{8d}{\pi^2},\qquad
 c=c_1^2+\kappa.
\end{equation}

\begin{theorem}[Sharp corner expansion]\label{thm:corner-expansion}
For every integer $n\geq2$,
\begin{align}
 \Lambda_c(n)&=A(N)^2+\kappa-\frac1N
                  +O\!\left(\frac{\log N}{N^2}\right)
                  \label{eq:asym-sharp-square}\\
 &=\frac4{\pi^2}\log^2N+b\log N+c-\frac1N
                  +O\!\left(\frac{\log N}{N^2}\right).
                  \label{eq:asym-sharp-expanded}
\end{align}
The implied constant is independent of $n$.  In particular the same
formula holds for both parities, and
\begin{equation}\label{eq:asym-reciprocal-limit}
 N\bigl(\Lambda_c(n)-A(N)^2-\kappa\bigr)\longrightarrow-1.
\end{equation}
\end{theorem}

This theorem concerns the actual corner value.  Its proof uses no
assertion about the location of the global maximum.

\subsection{Why a uniform row expansion is needed}

For $r+s=m$ and $x=r/m$, the harmonic row in
\eqref{eq:corner-harmonic-row} contains arguments close to both ends
of the interval.  Expanding it with $x$ fixed would miss rows with
$r$ or $s$ small.  Those rows cannot simply be discarded: the weight
$w(r/m)$ is comparable to $m/r$ at the left endpoint.

The remedy is to center at the exact integers and half integers first.
The uniform centering estimate proved in
Appendix~\ref{sec:asym-details-centering} reads
\begin{align}
 D_{m,r}={}&\psi(r+\tfrac12)+\psi(s)-\psi((1-x)/2)-\psi(m)
      \notag\\[-2pt]
 &-\frac{r\psi'(r+\tfrac12)+s\psi'(s)-s\psi'(m)}{2m}
      +E_{m,r},
 \qquad |E_{m,r}|\leq\frac{5}{4m^2}.
 \label{eq:asym-uniform-centering}
\end{align}
Since $\sum_{r<m}w(r/m)=O(m\log m)$, the error after weighting
and normalization is $O(\log m/m^2)$, at the scale required by the
theorem.  This treats $r=1$ and $s=1$ with the same bound.

One small-index term survives this procedure.  If
$u_r=r\psi'(r+1/2)-1$, then $u_r=O(r^{-2})$, but its contribution
to the coefficient of $1/m$ is nonzero:
\begin{equation}\label{eq:asym-small-index-sum}
 \sum_{r=1}^\infty\frac{u_r}{r}
       =1+2\log2-\frac{\pi^2}{4}.
\end{equation}
Keeping this convergent endpoint sum is as important as controlling
the large indices.

\subsection{The two contributions through first order}

All the integral quantities in the next proposition have removable
endpoint singularities.  With $w$ and $B$ from
\eqref{eq:corner-w-B}, put
\begin{equation}\label{eq:asym-integrals}
 I_1=\int_0^1\left(\frac{w(x)}{1-x}-\frac ax\right)dx,
 \qquad I_2=\int_0^1xw(x)\,dx,
 \qquad I_B=\int_0^1xw(x)B(x)\,dx.
\end{equation}

\begin{proposition}[Harmonic and smooth expansions]
\label{prop:asym-two-expansions}
There are real constants $C_H,C_B$, independent of the degree, such
that, for $m=n+1\geq3$,
\begin{align}
 \mathcal H_m
   &=A(m)^2+C_H+\frac{-a^2\log m+\beta_H}{m}
                    +O\!\left(\frac{\log m}{m^2}\right),
                    \label{eq:asym-harmonic-expansion}\\
 \beta_H&=-2-a^2(\gamma+2\log2)-a(I_1+I_2),
                    \label{eq:asym-beta-H}\\
 \frac{4V_n}{nm}
   &=C_B+\frac{\beta_B}{n}+R_{B,n},\qquad
          |R_{B,n}|\leq\frac{17}{n^2},
                    \label{eq:asym-smooth-expansion}\\
 C_B&=4\int_0^1w(t)\int_0^tB(x)\,dx\,dt,
 \qquad \beta_B=-a^2+2I_B+aI_2.
                    \label{eq:asym-beta-B}
\end{align}
In fact $-a^2<\beta_B<0$.
\end{proposition}

The proof is in Appendices~\ref{sec:asym-details-harmonic} and
\ref{sec:asym-details-smooth}.  Two features explain its structure.
For the harmonic part, subtracting the elementary singularity
$a(\log x+\gamma+2\log2)/x$ leaves a function whose second
derivative is bounded by a constant times
$x^{-1}+(1-x)^{-1}$.  Estimating its two endpoint cells directly
then gives the logarithmic quadrature error in
\eqref{eq:asym-harmonic-expansion}.  For the smooth part, the samples
in row $r$ are the midpoints of cells whose union is exactly
\[
 \left[-\frac{t}{2n},\ t+\frac{t}{2n}\right],\qquad t=\frac rm.
\]
The two short extensions of $[0,t]$ produce a first-order term.
Thus smoothness of $B$ does not make its $1/n$ correction vanish.

\subsection{The cancellation and the half-integer scale}

\begin{proof}[Proof of Theorem~\ref{thm:corner-expansion}]
The decisive identity among the remaining integrals is
\begin{equation}\label{eq:asym-integral-cancellation}
 2I_B-aI_1=a^2-1+a^2\log(\pi/2).
\end{equation}
It is useful to verify this before evaluating any integrals separately.
The double-angle identity
$\cot(\pi x)=(w(x)-w(x)^{-1})/2$ shows that the combined
integrand on the left is
\[
 x(w(x)^2-1)-2aw(x)+\frac{a^2}{x}.
\]
Since $w'=-(1+w^2)/a$, its integral over $[\epsilon,1]$ equals
\[
 a\epsilon w(\epsilon)-1+\epsilon^2
       +a^2\log\frac{\sin(\pi\epsilon/2)}{\epsilon}.
\]
Letting $\epsilon\downarrow0$ proves
\eqref{eq:asym-integral-cancellation}, including the endpoint term
$a^2-1$.

The coefficient $\beta_B$ is a constant, so replacing $\beta_B/n$
by $\beta_B/m$ in \eqref{eq:asym-smooth-expansion} costs only
$O(m^{-2})$.  Insert both expansions into the exact identity
\eqref{eq:corner-exact-decomposition}.  Its $2/m$ must be retained.
The coefficient of $1/m$, apart from $-a^2\log m/m$, becomes
\begin{align*}
 2+\beta_H+\beta_B
 &=-a^2-a^2(\gamma+2\log2)+2I_B-aI_1\\
 &=-1-a^2d.
\end{align*}
Consequently, for a real constant $C=-1+C_H+C_B$,
\begin{equation}\label{eq:asym-m-expansion}
 \Lambda_c(n)=A(m)^2+C
       -\frac{a^2\log m+1+a^2d}{m}
       +O\!\left(\frac{\log m}{m^2}\right).
\end{equation}
Lemmas~\ref{lem:corner-one-dimensional} and
\ref{lem:corner-defect-limit} identify $C=\kappa$: indeed
$L_n^2-A(n)^2=O(\log n/n^2)$ and
$A(m)^2-A(n)^2\to0$.

Finally the half step $N=m-1/2$ absorbs the logarithmic part of the
first correction.  Taylor's formula, with a uniform bound on the second
derivative of $A(z)^2$, gives
\[
 A(N)^2=A(m)^2-\frac{a^2\log m+a^2d}{m}
                       +O\!\left(\frac{\log m}{m^2}\right),
 \qquad \frac1N=\frac1m+O(m^{-2}).
\]
Substitution into \eqref{eq:asym-m-expansion} proves
\eqref{eq:asym-sharp-square}.  Expanding the square gives
\eqref{eq:asym-sharp-expanded}, with the constants
\eqref{eq:asym-coefficients}, and multiplication by $N$ gives
\eqref{eq:asym-reciprocal-limit}.  All estimates used above hold
on the full integer sequence with constants independent of the degree.
\end{proof}
